% GENERATED FILE. Edit canonical lessons, then run npm run generate:books.
% canonical-source-hash: fnv1a64:e93448353f0564e7
% canonical-lessons: FA-C410-hezb, FA-C410-layehe, FA-C410-tasvib-kardan, FA-C410-sanad, FA-C410-tabeiyyat

\chapter{Parties and new laws}
\label{ch:fa-parties-and-new-laws}
\hlchaptermodality{410}

\begin{chapteropening}
\textbf{By the end of this chapter:} \emph{I can talk about political parties, say that a bill was sent to parliament and passed, and talk about documents and citizenship.}

Complete the last lesson of chapter 410: I can talk about political parties, say that a bill was sent to parliament and passed, and talk about documents and citizenship.
\end{chapteropening}

\section[\texorpdfstring{\fa{حزب} (hezb)}{حزب (hezb)}]{\fa{حزب} (hezb) — a political party}

\label{lesson:FA-C410-hezb}



\begin{quote}
Before the new one: say the Persian for to vote, then the Persian for parliament.
\end{quote}

\subsection*{You'll want to know: \fa{حزب}}
\textbf{\fa{حزب}} (\emph{hezb}) — \textquotedblleft{}a political party\textquotedblright{}.

\begin{quote}
\fa{او} \fa{عضو} \fa{یک} \fa{حزب} \fa{کوچک} \fa{است}.
\end{quote}

\begin{quote}
\emph{u 'ozv-e yek hezb-e kuchek ast.}
\end{quote}

\begin{quote}
She is a member of a small party.
\end{quote}

Plural \emph{ahzâb}, an Arabic broken plural.

\begin{cousinweb}
\textbf{\fa{حزب}} was borrowed from Arabic \emph{ḥizb} \textquotedblleft{}a party, a group\textquotedblright{} (EtymDB).
\end{cousinweb}

\begin{grammarlens}[title={What you've built}]
One more word for telling, explaining and describing at B1.
\end{grammarlens}

\subsection*{Your turn}
\begin{itemize}
\raggedright
  \item \emph{Say it:} \emph{hezb}
  \item \emph{Say it:} \emph{hezb}, once more
  \item \emph{Say it:} a sentence of your own with \emph{hezb}
  \item \emph{Recall:} say the Persian for to vote, then the Persian for parliament, then say \emph{hezb} again
\end{itemize}

\begin{tcolorbox}[breakable,colback=teal!4,colframe=teal!35!black,title={Before you move on}]
What does \fa{حزب} mean? (\textquotedblleft{}A political party\textquotedblright{}.) Say it once more.
\end{tcolorbox}
\section[\texorpdfstring{\fa{لایحه} (lâyehe)}{لایحه (lâyehe)}]{\fa{لایحه} (lâyehe) — a draft law, a bill (in parliament)}

\label{lesson:FA-C410-layehe}



\begin{quote}
Before the new one: say the Persian for parliament, then the Persian for a political party.
\end{quote}

\subsection*{You'll want to know: \fa{لایحه}}
\textbf{\fa{لایحه}} (\emph{lâyehe}) — \textquotedblleft{}a draft law, a bill (in parliament)\textquotedblright{}.

\begin{quote}
\fa{دولت} \fa{یک} \fa{لایحه} \fa{تازه} \fa{به} \fa{مجلس} \fa{فرستاد}.
\end{quote}

\begin{quote}
\emph{dowlat yek lâyehe-ye tâze be majles ferestâd.}
\end{quote}

\begin{quote}
The government sent a new bill to parliament.
\end{quote}

A government bill; once passed it becomes a \emph{qânun}.

\begin{cousinweb}
\textbf{\fa{لایحه}} is usually traced to Arabic \emph{lā'iḥa} \textquotedblleft{}a written statement, a memorandum\textquotedblright{}.
\end{cousinweb}

\begin{grammarlens}[title={What you've built}]
One more word for telling, explaining and describing at B1.
\end{grammarlens}

\subsection*{Your turn}
\begin{itemize}
\raggedright
  \item \emph{Say it:} \emph{lâyehe}
  \item \emph{Say it:} \emph{lâyehe}, once more
  \item \emph{Say it:} a sentence of your own with \emph{lâyehe}
  \item \emph{Recall:} say the Persian for parliament, then the Persian for a political party, then say \emph{lâyehe} again
\end{itemize}

\begin{tcolorbox}[breakable,colback=teal!4,colframe=teal!35!black,title={Before you move on}]
What does \fa{لایحه} mean? (\textquotedblleft{}A draft law, a bill (in parliament)\textquotedblright{}.) Say it once more.
\end{tcolorbox}
\section[\texorpdfstring{\fa{تصویب} \fa{کردن} (tasvib kardan)}{تصویب کردن (tasvib kardan)}]{\fa{تصویب} \fa{کردن} (tasvib kardan) — to pass (a law), to approve}

\label{lesson:FA-C410-tasvib-kardan}



\begin{quote}
Before the new one: say the Persian for a political party, then the Persian for a draft law.
\end{quote}

\subsection*{You'll want to know: \fa{تصویب} \fa{کردن}}
\textbf{\fa{تصویب} \fa{کردن}} (\emph{tasvib kardan}) — \textquotedblleft{}to pass (a law), to approve\textquotedblright{}.

\begin{quote}
\fa{مجلس} \fa{این} \fa{لایحه} \fa{را} \fa{تصویب} \fa{کرد}.
\end{quote}

\begin{quote}
\emph{majles in lâyehe râ tasvib kard.}
\end{quote}

\begin{quote}
Parliament passed this bill.
\end{quote}

\emph{Tasvib} \textquotedblleft{}approval\textquotedblright{} + \emph{kardan}. \emph{Tasvib shodan} — to be passed.

\begin{cousinweb}
\textbf{\fa{تصویب}} is usually traced to Arabic \emph{taṣwīb} \textquotedblleft{}approving as right\textquotedblright{}; with \emph{kardan} \textquotedblleft{}to do\textquotedblright{} it means to pass a law.
\end{cousinweb}

\begin{grammarlens}[title={What you've built}]
One more word for telling, explaining and describing at B1.
\end{grammarlens}

\subsection*{Your turn}
\begin{itemize}
\raggedright
  \item \emph{Say it:} \emph{tasvib kardan}
  \item \emph{Say it:} \emph{tasvib kardan}, once more
  \item \emph{Say it:} a sentence of your own with \emph{tasvib kardan}
  \item \emph{Recall:} say the Persian for a political party, then the Persian for a draft law, then say \emph{tasvib kardan} again
\end{itemize}

\begin{tcolorbox}[breakable,colback=teal!4,colframe=teal!35!black,title={Before you move on}]
What does \fa{تصویب} \fa{کردن} mean? (\textquotedblleft{}To pass (a law), to approve\textquotedblright{}.) Say it once more.
\end{tcolorbox}
\section[\texorpdfstring{\fa{سند} (sanad)}{سند (sanad)}]{\fa{سند} (sanad) — a document; a deed}

\label{lesson:FA-C410-sanad}



\begin{quote}
Before the new one: say the Persian for a draft law, then the Persian for to pass.
\end{quote}

\subsection*{You'll want to know: \fa{سند}}
\textbf{\fa{سند}} (\emph{sanad}) — \textquotedblleft{}a document; a deed\textquotedblright{}.

\begin{quote}
\fa{سند} \fa{خانه} \fa{به} \fa{نام} \fa{مادرم} \fa{است}.
\end{quote}

\begin{quote}
\emph{sanad-e khâne be nâm-e mâdaram ast.}
\end{quote}

\begin{quote}
The deed of the house is in my mother's name.
\end{quote}

Plural \emph{asnâd}, an Arabic broken plural.

\begin{cousinweb}
\textbf{\fa{سند}} was borrowed from Arabic \emph{sanad} (EtymDB), a word usually glossed as a support, what a claim leans on; in Persian it is a document or a deed.
\end{cousinweb}

\begin{grammarlens}[title={What you've built}]
One more word for telling, explaining and describing at B1.
\end{grammarlens}

\subsection*{Your turn}
\begin{itemize}
\raggedright
  \item \emph{Say it:} \emph{sanad}
  \item \emph{Say it:} \emph{sanad}, once more
  \item \emph{Say it:} a sentence of your own with \emph{sanad}
  \item \emph{Recall:} say the Persian for a draft law, then the Persian for to pass, then say \emph{sanad} again
\end{itemize}

\begin{tcolorbox}[breakable,colback=teal!4,colframe=teal!35!black,title={Before you move on}]
What does \fa{سند} mean? (\textquotedblleft{}A document; a deed\textquotedblright{}.) Say it once more.
\end{tcolorbox}
\section[\texorpdfstring{\fa{تابعیت} (tâbe'iyat)}{تابعیت (tâbe'iyat)}]{\fa{تابعیت} (tâbe'iyat) — citizenship, nationality}

\label{lesson:FA-C410-tabeiyyat}



\begin{quote}
Before the new one: say the Persian for to pass, then the Persian for a document.
\end{quote}

\subsection*{You'll want to know: \fa{تابعیت}}
\textbf{\fa{تابعیت}} (\emph{tâbe'iyat}) — \textquotedblleft{}citizenship, nationality\textquotedblright{}.

\begin{quote}
\fa{او} \fa{تابعیت} \fa{دو} \fa{کشور} \fa{را} \fa{دارد}.
\end{quote}

\begin{quote}
\emph{u tâbe'iyat-e do keshvar râ dârad.}
\end{quote}

\begin{quote}
He has the citizenship of two countries.
\end{quote}

\emph{Tâbe'iyyat gereftan} — to become a citizen.

\begin{cousinweb}
\textbf{\fa{تابعیت}} is usually traced to Arabic \emph{tābi'iyya}, built on \emph{tābi'} \textquotedblleft{}following, subject to\textquotedblright{}: a citizen belongs to a state.
\end{cousinweb}

\begin{grammarlens}[title={What you've built}]
One more word for telling, explaining and describing at B1.
\end{grammarlens}

\subsection*{Your turn}
\begin{itemize}
\raggedright
  \item \emph{Say it:} \emph{tâbe'iyat}
  \item \emph{Say it:} \emph{tâbe'iyat}, once more
  \item \emph{Say it:} a sentence of your own with \emph{tâbe'iyat}
  \item \emph{Recall:} say the Persian for to pass, then the Persian for a document, then say \emph{tâbe'iyat} again
\end{itemize}

\begin{tcolorbox}[breakable,colback=teal!4,colframe=teal!35!black,title={Before you move on}]
What does \fa{تابعیت} mean? (\textquotedblleft{}Citizenship, nationality\textquotedblright{}.) Say it once more.
\end{tcolorbox}
